\section{Introduction}
\label{sec:introduction}

Knowledge graphs (KGs) connect facts across sources and support question
answering, recommendation, and data analysis \citep{ref_ji2021,ref_pan2023}.
These applications depend on the quality of the stored facts.
A missing relation can hide useful information, while a wrong value can
lead to a wrong answer or decision.
As more graphs are built from text, checking and improving their quality
has become an important part of KG construction
\citep{ref_zaveri2016,ref_paulheim2017}.
The goal is to recover useful facts and preserve the correct information
already in the graph.

Text-based extraction leaves several kinds of errors in the same graph.
Source noise and unclear wording can cause missing fields, repeated facts,
wrong relations, and values assigned to the wrong fields
\citep{ref_ji2021,ref_paulheim2017}.
These errors connect graph structure with the meaning of the source text.
For example, both a subtotal and a final amount may appear on a receipt,
but they serve different roles.
Checking whether an amount appears in the text is one part of deciding
whether it belongs in the total field.
Graph repair therefore needs both structural checks and evidence for each
field value.

Existing methods provide several ways to improve graph quality.
Rule-based checks find violations of graph and field requirements
\citep{ref_kontokostas2014,ref_shacl2017}.
Embedding models score how likely a fact is from patterns in the graph
\citep{ref_bordes2013,ref_nickel2016}.
Large language models (LLMs) use text and graph context to suggest repairs
\citep{ref_pan2023,ref_lin2025}.
Each approach provides useful information for a repair decision.
A rule violation points to a broken requirement.
A model proposal offers a possible replacement.
Connecting these signals requires a common view of the errors, clear
bounds on changes, and a way to choose among repair actions.

Quality errors arise at different scales.
Local checks can find repeated facts, while graph checks test relations
and conflicts across fields.
Source checks connect stored values to the document evidence
\citep{ref_zaveri2016,ref_wang2021}.
A value may satisfy a field rule yet refer to the wrong part of a document.
This makes it important to bring the checks together when planning a repair.
We develop \emph{Multi-Scale Quality Assessment} across entity, graph,
and context scales.
The checks form a shared assessment profile that records errors and
quality scores at each scope.
This profile links the detected problems to the repair process and
provides a common basis for assessing proposed changes.

Repair also involves a trade-off between recovering missing facts and
preserving correct ones.
Adding a fact can create a field conflict, while deleting a fact can
remove useful information.
We introduce \emph{Constraint-Driven Optimization} to give these changes
clear bounds.
Source and field constraints define which generated triples can enter
the repaired graph.
For edit selection, quality lower bounds and a density limit define
allowed trial graphs, while edit costs favor smaller changes.
These constraints connect the repair goal with explicit requirements
for the output.
They also make the grounds for accepting or rejecting a proposed change
clear.

The choice of repair action must also reflect the type of error.
Simple rules can remove repeated facts and correct clear structural errors.
Missing or wrong values require reading the source text.
Several proposed edits may target the same field or change the effect
of another edit.
We develop \emph{Scale-Aware Routing} to guide this choice using the
assessment profile.
A small neural model assigns weights to edits at the entity, graph,
and context scales.
The selector combines these weights with quality gains and edit costs,
then compares the proposed changes under the graph bounds.
Reassessment after an accepted edit guides the next choice.
This connects the scope of the detected problem with the scope and cost
of the repair.

Together, these designs link assessment, repair constraints, and action
selection in one framework.
We study document graphs whose fields must follow a given schema and
use values from the available source text.
The main path combines rule-based cleanup, LLM generation, and output
filtering.
An error report or field-specific source-line hints guide generation.
The routing path selects edits one at a time and checks the graph after
each accepted change.
We evaluate the generation path and the routing path with controls for
the components in each path.
This setup shows how the proposed designs affect the graph at each
stage of repair.

The experiments cover three Chinese record collections, SROIE receipts,
and CUAD contracts, together with human review of graph differences
and system edits.
On 225 graphs with 450 added defects, Diagnosis + Gate repairs 98.00\%
of the defects.
On model-built graphs, it raises triple F1 from 87.99\% to 92.48\%.
On a new test of 60 receipts, source-line hints raise F1 over Simple
Pipeline by 4.76 points.
Shared-prompt comparisons and component tests examine the roles of
generation, filtering, evidence context, and edit selection.
The evaluation tracks recovered fields, retained facts, request cost,
and human acceptance of edits.

The paper makes three contributions:
\begin{itemize}
    \item \textbf{Multi-scale quality assessment.}
    We combine local structure, graph rules, and source checks in a shared
    profile that connects detected defects with repair decisions.
    \item \textbf{Constraint-driven repair and scale-aware routing.}
    We define source and field constraints for generated graphs.
    A neural routing path weighs candidate edits and compares their
    quality gains and costs under graph bounds.
    \item \textbf{Evaluation on document graphs.}
    We compare repair methods and their components on added defects,
    model outputs, receipts, and contracts, and include human review.
    Saved outputs and decision logs support checks of each repair stage.
\end{itemize}

Section~\ref{sec:related_work} reviews related work.
Sections~\ref{sec:framework} and~\ref{sec:implementation} present the
method and its implementation.
Section~\ref{sec:experiments} reports the experiments and results.
Section~\ref{sec:conclusion} concludes the paper.
